%%%%%%%%%%%%%%%%%%%%%%%%%%%%%%%%%%%%%%%%%
% Important note:
% This template requires the resume.cls file to be in the same directory as the
% .tex file. The resume.cls file provides the resume style used for structuring the
% document.
%
%%%%%%%%%%%%%%%%%%%%%%%%%%%%%%%%%%%%%%%%%

%----------------------------------------------------------------------------------------
%	PACKAGES AND OTHER DOCUMENT CONFIGURATIONS
%----------------------------------------------------------------------------------------

\documentclass{resume} % Use the custom resume.cls style

\usepackage[left=0.75in,top=0.6in,right=0.75in,bottom=0.6in]{geometry} % Document margins
\usepackage[hyperref]{}
\usepackage{enumitem}
\usepackage[colorlinks=true, urlcolor=blue, linkcolor=red]{hyperref}
\usepackage{comment}
\usepackage{fancyhdr} 
\pagestyle{fancy}
%\fancyhf{}
\renewcommand{\headrulewidth}{0pt}

%\fancyhf{}
%\cfoot{\thepage}
%\pagestyle{fancy} 

\newcommand{\tab}[1]{\hspace{.2667\textwidth}\rlap{#1}}
\newcommand{\itab}[1]{\hspace{0em}\rlap{#1}}
\name{\large{Suryanarayana Maddu Kondaiah}}\\ % Your name

\address{\small \textit{Center for Computational Biology, Flatiron Institute, 162 5th Ave, New York, NY 10010\vspace{-0.5em}}} 
\address{\small \textit{Department of Molecular and Cellular Biology, Harvard University, 52 Oxford St, Cambridge, MA 02138} } 
\address{\small \href{smaddu@flatironinstitute.org}{smaddu@flatironinstitute.org} \vert \href{https://scholar.google.de/citations?user=oL3QYLkAAAAJ&hl=en}{\:  \textrm{Google Scholar }  } \vert \href{https://orcid.org/0000-0003-2702-4427}{ \textrm{ ORCID: 0000-0003-2702-4427}\vspace{-1.5em}}}
% Your address
%\address{123 Pleasant Lane \\ City, State 12345} % Your secondary addess (optional)
%\address{(+country code) phone number\\ emailhere@email.com} % Your phone number and email

\begin{document}

%----------------------------------------------------------------------------------------
%	EDUCATION SECTION
%----------------------------------------------------------------------------------------

\begin{rSection}{Research Interests}
Scientific Machine Learning, Biophysical Modeling, Statistical Inference, Computational Physics
% My research lies at the interface of physical modeling, scientific computing, and AI-assisted biophysics. 
% I develop scalable AI methods for inferring stochastic processes and dynamical systems from high-dimensional, noisy data, 
% with emphasis on flow matching, score-based inference, and physics-informed neural networks. 
% I am particularly interested in integrating open-source AI tooling with mechanistic modeling to enable 
% scientific discovery in biology and the physical sciences.
\end{rSection}

\begin{rSection}{\small{Appointments }}
{\small Research Associate, Department of Molecular
and Cellular Biology} \hfill {\small \textit{ Dec 2025 - present}} \\
{\small \bf Harvard University, Cambridge, MA} \\
{\small \textit{Advisor:}  Prof. Daniel J. Needleman} \\ \\
{\small CCB$_X$ Fellow, Center for Computational Biology } \hfill {\small \textit{Dec 2025 - present}} \\
{\small Flatiron Research Fellow, Center for Computational Biology  } \hfill {\small \textit{March 2022 - Nov 2025}} \\
{\small \bf Flatiron Institute, Simons Foundation, NY } \\
{\small \textit{Advisor:} Prof. Michael J. Shelley (Flatiron Institute/Courant Institute) }
\end{rSection}


\begin{rSection}{\small{Education}}
%--copy and paste this region  if you need more--
%{\small \bf Technical University of Dresden, Germany \& } \\
{\small \bf Max Planck Institute of Molecular Cell Biology and Genetics \& } \\
{\small \bf Technische Universit{\"a}t Dresden, Dresden, Germany}
\hfill {\small \textit{2017 - 2021} }
\\ \small Doctor of Philosophy (Ph.D.) in Computer Science, \textit{summa cum laude}
\\ \small \textit{Thesis advisor(s)}: Prof. Ivo F. Sbalzarini \& Prof. Christian L. M{\"u}ller  %\hfill { GPA: your gpa }
% \\ \small \textit{Thesis title}: Data-driven modeling and simulation of spatiotemporal processes with a view toward applications in biology
\\ \\
{\small \bf R{\"u}hr-Universit{\"a}t Bochum, Bochum, Germany} \hfill {\textit{2014 - 2016} }
\\ \small Master of Science (M.Sc) in Computational Science and Engineering
\\ \\
{\small \bf National Institute of Technology, Karnataka, India} \hfill {\textit{2010 - 2014} }
\\ \small Bachelor of Technology (B.Tech) in Mechanical Engineering
%--copy and paste this region  if you need more--

\end{rSection}

\begin{comment}
\begin{rSection}{\small Research Experience}

{\small \bf Center for Computational Biology, Flatiron Institute, New York, USA} \hfill {\small \textit{2022 - present} }\\
{\small \textit{Flatiron Research Fellow} 


{\small \bf Max Planck Institute of Molecular Cell Biology and Genetics, Dresden, Germany} \hfill {\small \textit{2017 - 2022}} \\
{\small \textit{Doctoral Researcher} 

{\small \bf Center for Computational Mathematics, Flatiron Institute, New York, USA} \hfill {\small \textit{Summer 2019}} \\
{\small \textit{Intern} 
\vspace{-0.5em}


\end{rSection}
\end{comment}

\begin{rSection}{\small Selected Publications}
%--copy and paste this region  if you need more--

\begin{enumerate}[leftmargin=*]
    \item \small \textbf{Maddu, S.}, Chard\`es, V., Shelley, M. Learning biophysical models of gene regulatory dynamics using probability flow matching. \textit{Submitted (2026)} 
    \item \small \textbf{Maddu, S.$^{\dag}$}, Kelleher, C.P$^{\dag}$.,  Basaran, M., Müller-Reichert, T., Shelley, M.J., Needleman, D.J. Active liquid crystal theory explains the collective organization of microtubules in human mitotic spindles. \textit{Proceedings of the National Academy of Sciences (2026)}
    \item \small Zhang, S., \textbf{Maddu, S.,} Qiu, X.,  Chard\`es, V. Inferring stochastic dynamics with growth from cross-sectional data. \textit{Neural Information Processing Systems (2025)}
    \item \small \textbf{Maddu, S.,} Chard\`es, V., Shelley, M. Inferring biological processes with intrinsic noise from cross-sectional data. \textit{Proceedings of the National Academy of Sciences (2025) }
    \item \small The Well: a Large-Scale Collection of Diverse Physics Simulations for Machine Learning. $\ddagger$ \textit{Datasets and Benchmarks Track, Neural Information Processing Systems (2024)}
    %\textit{arXiv:2410.07501}  - 
    \item \small Sturm, D., \textbf{Maddu, S.,} Sbalzarini, I.F. Learning locally dominant force balances in active particle systems. \textit{Proceedings of the Royal Society A (2024)}
    \item \small {\textbf{Maddu, S., et al.} Learning fast, accurate, and stable closures of a kinetic theory of an active fluid. \textit{Journal of Computational Physics, p.112869 (2024)}
    \item \small {Chard\`es, V\:$^\dag$., \textbf{Maddu, S\:$^\dag$.,} Shelley, M. Stochastic force inference via density estimation. \textit{AI for Scientific Discovery: From Theory to Practice workshop, Neural Information Processing Systems (2023)}
    % \item \small \textbf{Maddu, S., et al.} Learning deterministic hydrodynamic equations from stochastic active particle dynamics. \textit{arXiv:2201.08623} (2023) - under review at PRR
%\cventry{2022}{Learning deterministic hydrodynamic equations from stochastic active particle dynamics}{}{}{}{\textbf{Suryanarayana Maddu}, Quentin Vagne, I.F. Sbalzarini - under review}
    \item \small Perez, S., \textbf{Maddu, S.,} Sbalzarini, I.F., Poncet, P. Adaptive weighting of Bayesian physics informed neural networks for multitask and multiscale forward and inverse problems. \textit{Journal of Computational Physics, p.112342 (2023)} 
    \item \small \textbf{Maddu, S., et al.} STENCIL-NET for equation-free forecasting from data. \textit{Nature Scientific Reports, 13, 12787 (2023)} 
    \item \small \textbf{Maddu, S., et al.} Stability selection enables robust learning of partial differential equations from limited noisy data. \textit{Proceedings of the Royal Society A, 478(2262), p.20210916 (2022)} 
    \item \small Jonsson, J., Cheeseman, B., \textbf{Maddu, S.,} Gonciarz, K., Sbalzarini, I.F. Parallel discrete convolutions on adaptive particle representations of images. \textit{IEEE Transactions on Image Processing, 31, pp.4197-4212 (2022)} 
    \item \small \textbf{Maddu, S., et al.} Inverse-Dirichlet weighting enables reliable training of physics-informed neural networks. \textit{Machine Learning: Science and Technology, 3(1), p.015026 (2021)} 
    \item \small \textbf{Maddu, S., et al.} Learning physically consistent differential equation models from data using group sparsity. \textit{Physical Review E, 103(4), p.042310 (2021)} 
\end{enumerate}
$\dag$ \textit{Equal contribution} \: $\ddagger$ \textit{Domain expert}

\end{rSection}


\begin{rSection}{\small SELECTED INVITED AND CONTRIBUTED TALKS AND SEMINARS}
{\small \textit{Note: Poster presentations and summer schools not included}}
%--copy and paste this region  if you need more--
% { \small \textbf{Invited Talk,} D$^3$ International colloquium, Pavia, Italy} \hfill{\small \textit{October 2025}} \\
% {\small \textit{Learning accurate closures of a kinetic theory of an active fluid} }

{\small \textbf{Invited Talk,} Theory of Living Matter Seminar Series} \hfill{\small \textit{February 2026}} \\
\textit{Active Liquid Crystal Theory of Microtubule Organization in Human Spindles} \\
{\small \textbf{Invited Talk,} Soft Matter Seminar, Department of Physics, Syracuse University,  Syracuse} \hfill{\small \textit{February 2026}}\\
{\small \textit{Inference of Multiscale Dynamical Systems from Biological Data} }
\\
{\small \textbf{Invited Talk,} AMS special seminar series, Johns Hopkins University, Baltimore} \hfill{\small \textit{January 2026}}\\
{\small \textit{Inference of Multiscale Dynamical Systems from Biological Data} }
\\
{ \small \textbf{Invited Talk,} Physics of Life symposium, CZ Biohub SF/UCSF, San Francisco, California} \hfill{\small \textit{September 2025}}\\
{\small \textit{Active Liquid Crystal Theory of Microtubule Organization in Human Spindles} }
\\
{ \small \textbf{Invited Talk,} AI in Molecular Biology, Sante Fe, New Mexico} \hfill{\small \textit{September 2025}}\\
{\small \textit{Scalable Inference of Biophysical Models from Multi-Omics Data} }
\\
{ \small \textbf{Invited Talk,} Data-Driven Design of Resilient Meta Materials ($\textrm{D}^3$) colloquium, Pavia, Italy} \hfill{\small \textit{September 2025}} 
{\small \textit{Data-driven closure modeling of interacting particle assemblies} }
\\
% { \small \textbf{Invited Talk,} Frontiers in Applied and Computational Mathematics, NJIT, New Jersey} \hfill{\small \textit{June 2025}} \\
% {\small \textit{Inference of biophysical models from cross-sectional datasets} }
{ \small \textbf{Invited Talk,} Frontiers in Applied \& Computational Mathematics, NJIT, New Jersey} \hfill{\small \textit{June 2025}} \\
{\small \textit{Scalable inference of biological processes from cross-sectional datasets} }
\\
{ \small \textbf{Invited Talk,} Workshop on Modeling and Inference of Stochastic Processes, Flatiron Institute} \hfill{\small \textit{November 2024}} \\
{\small \textit{Inferring biological processes with intrinsic noise from cross-sectional data} }
\\
{ \small \textbf{Invited Talk,} Flatiron-wide Algorithms and Mathematics, New York} \hfill{\small \textit{October 2024}} \\
{\small \textit{Integrating mechanistic models with single-cell data via probability flow inference} }
\\
{  \small \textbf{Invited Talk,} \small X International Conference on Computational Bioengineering (ICCB), Vienna} \hfill {\small \textit{September 2023}}\\
{\small \textit{Probing the physical basis of human mitotic spindle self-organization using data from electron
microscopy}}
\\
{ \small \textbf{Invited Talk,} {APS March Meeting, Las Vegas}}  \hfill {\small \textit{March 2023}} \\
{\small \textit{Learning accurate closures of a kinetic theory of an active fluid} }
\\
{\small \textbf{Talk,} Conference on Parallel Processing for Scientific Computing (SIAM-PP), Seattle} \hfill {\small \textit{February 2022}}\\
{\small \textit{Learning solution adaptive discretization stencils from data}}
% {\small \textbf{Seminar,} Harvard Applied Science Seminar} \hfill {\small \textit{October 2021}}\\
% {\small \textit{Dynamical systems inference from data with a view toward applications in biology}}
\\
{\small \textbf{Talk,}  World Congress in Computational Mechanics and ECCOMAS Congress, Paris} \hfill {\small \textit{January 2021}}\\
{\small \textit{Encoding knowledge with group sparsity for model learning from limited and noisy biological data}}
% \\
% {\small \textbf{Talk,}  World Congress in Computational Mechanics and ECCOMAS Congress, Paris} \hfill {\small \textit{January 2021}}\\
% {\small \textit{Learning computable models from data}}
\\
{ \small \textbf{Invited Talk,} New York Scientific Data Summit (NYSDS), Columbia University, New York} \hfill{\small \textit{June 2019}} \\
{\small \textit{Enabling data-driven modeling in biology by statistical learning of interpretable mathematical models from microscopy data} }
\\
{\small \textbf{Invited talk,} Center for Computational Mathematics, Flatiron Institute, New York} \hfill {\small \textit{June 2019}}\\
{\small \textit{Enabling data-driven modeling in biology}}

%\begin{rSection}{\small Poster presentation}




%--copy and paste this region  if you need more--
\end{rSection}
%----------------------------------------------------------------------------------------
%	SKILLS SECTION
%----------------------------------------------------------------------------------------
\newpage
\begin{rSection}{\small Teaching Experience}
{\small \bf Teaching Assistant} \hfill{ \small \textit{Fall 2017 }} \\
{\small Instructed 50+ students in Prof. Ivo F. Sbalzarini’s course on \textit{Particle Methods and Computer Graphics} at the Computer Science Department, TU Dresden}
\\
{\small \bf Teaching Assistant}} \hfill{ \small \textit{Fall 2018 }} \\
{\small Assisted in teaching and developing course materials, including exercises and exams, for Prof. Ivo F. Sbalzarini’s \textit{Basic Numerical Methods} course in the Computational Modeling and Simulation Master's program at TU Dresden.}
\\
{\small \textbf{Tutor} at \textit{Dresden Summer School in Systems Biology, Dresden}} \hfill{ \small \textit{August 2018, 2019 } }
\\
{\small \textbf{Invited Lecturer} at \textit{ScaDS Big Data summer school, Dresden}} \hfill{ \small \textit{August 2019 } }
\\
{\small \textbf{Invited Lecturer} at \textit{Biophysics Summer School, Crete}} \hfill{ \small \textit{August 2025 } }

\end{rSection}

\begin{rSection}{Software and Open-Source} 
{\small \textbf{\href{https://git.mpi-cbg.de/mosaic/software/machine-learning/pde-stride}{PDE-STRIDE}}: Stability-inference framework for robust identification of differential equations from data.} 
\\
{\small \textbf{\href{https://git.mpi-cbg.de/mosaic/software/machine-learning/inverse-dirichlet-pinn}{Inverse-Dirichlet PINN}}: Reliable training of PINNs for multi-scale dynamical systems.  } 
\\
{\small \textbf{\href{https://git.mpi-cbg.de/mosaic/software/machine-learning/stencil-net}{STENCIL-NET}}: Equation-free forecasting from data using learned discretization stencils.} 
\\
{\small \textbf{\href{https://polymathic-ai.org/the_well/}{The WELL Dataset}}: Large-scale collection of diverse physics simulations for machine learning (NeurIPS 2024 Datasets and Benchmarks).} \\
{\small \textbf{\href{https://github.com/vchz/pfi}{PFI}}: Stochastic Force Inference}

\end{rSection}




\begin{rSection}{\small Mentoring Experience}
{\small \textbf{Master thesis supervision}:  \textit{Happy Das (TU Dresden, 2020); Dominik Sturm (TU Dresden, 2021)  }} \\
{\small \textbf{Interns supervision}:  \textit{Koushiki Dasgupta Chaudhuri (IIT Madras, 2021); Nirbhay Patil (IIT Bombay, 2018); Th\`eo Uscidda (ENSAE, 2024) }} \\
{\small \textbf{Doctoral thesis supervision}: \textit{Sarah Perez (University of Pau, 2022) along with Prof. Philip Poncet}; \textit{Max Rosenkranz (DFG graduate school on data-driven modeling of materials, TU Dresden, 2023 - present)} 
\end{rSection}


\begin{rSection}{Awards and Honors} 
% {\small \textbf{Otto Hahn award nomination} } \hfill{\small \textit{2022}}\\
% % {\small One of three doctoral theses nominated by the Max Planck Institute for Molecular Cell Biology \& Genetics}

{\small \textbf{Simons Foundation CCBx fellowship} } \hfill{\small \textit{2025}}
% {\small Research grant from Simons Foundation}
\\
{\small \textbf{Harvard QBio initiative visiting scholar grant} } \hfill{\small \textit{Summer 2022}}
% {\small Recipient of research grant at the NSF-Simons Center for Mathematical & Statistical Analysis of Biology, Harvard University}
\\
{\small \textbf{Karnataka State Government Scholarship for Higher Education Abroad} } \hfill{\small \textit{2014 - 2016}}
% {\small Recipient of the prestigious Karnataka State Government Scholarship for Higher Education Abroad, which supported my postgraduate studies in Germany}

%{One of the few students selected for the International Guest Student program on Scientific Computing at J{\"u}lich Supercomputing Center and IBM, Germany 2016.}

%{\small Distinctive performance in 6th Indian National Science Olympiad, 5th and 6th International level Science Talent Examination}

\end{rSection}


\begin{rSection}{Professional service} 
{\small \textbf{Reviewing:} Nature Communications, Proceedings of the National Academy of Sciences (PNAS), Proceedings of the Royal Society A, Journal of Computational Physics (JCP), Journal of Computational and Applied Mathematics (JCAM), Physical Review E (PRE), Physical Review X Life, Nature Communications Physics, Physical Review Letters (PRL)} 
\\
{\small \textbf{Reviewing duties for ML conferences:} NeurIPS 2024, ICML 2024, ICML 2025} 
\\
{\small \textbf{Area Chair:} AI for Science workshop in ICML 2024, AI for Science workshop NeurIPS, 2025}  
\\
{\small \textbf{Organizer:} Workshop on Modeling and Inference of Stochastic Processes in Cells, Flatiron Institute, New York, November 2024} 

\end{rSection}

\end{document}----------------------------

